\documentclass[10pt]{article}
\usepackage[T1]{fontenc}
\usepackage{newtxtext}
\usepackage{newtxmath}
\usepackage[letterpaper,top=0.68in,bottom=0.70in,left=0.72in,right=0.72in,
  headheight=12pt,headsep=7pt,footskip=20pt]{geometry}
\usepackage{amsmath,mathtools,bm}
\usepackage{booktabs,tabularx,array}
\usepackage{enumitem}
\usepackage{xcolor}
\usepackage{titlesec}
\usepackage{caption}
\usepackage{fancyhdr}
\usepackage[authoryear,round]{natbib}
\usepackage{hyperref}
\usepackage{url}

% Self-contained compact OR-Bench layout.
\definecolor{orblue}{RGB}{24,67,108}
\definecolor{orgray}{RGB}{88,94,101}
\definecolor{orlight}{RGB}{246,247,248}
\hypersetup{
  colorlinks=true,
  linkcolor=orblue,
  citecolor=orblue,
  urlcolor=orblue,
  pdfborder={0 0 0}
}
\setlength{\parindent}{1em}
\setlength{\parskip}{0pt}
\setlength{\tabcolsep}{4.5pt}
\renewcommand{\arraystretch}{1.08}
\setlist[itemize]{leftmargin=1.35em,itemsep=1pt,topsep=3pt,parsep=0pt}
\setlist[enumerate]{leftmargin=1.45em,itemsep=1pt,topsep=3pt,parsep=0pt}
\captionsetup{font=small,labelfont=bf,labelsep=period,skip=4pt,hypcap=false}
\titleformat{\section}
  {\bfseries\fontsize{11.6}{13.2}\selectfont}
  {\thesection.}{0.45em}{}
\titleformat{\subsection}
  {\bfseries\fontsize{10.2}{11.8}\selectfont}
  {\thesubsection.}{0.45em}{}
\titlespacing*{\section}{0pt}{10pt plus 2pt minus 2pt}{3pt}
\titlespacing*{\subsection}{0pt}{7pt plus 1pt minus 1pt}{2pt}

\makeatletter
\newcommand{\obshorttitle}{OR-Bench Candidate}
\newcommand{\shorttitle}[1]{\renewcommand{\obshorttitle}{#1}}
\AtBeginDocument{\hypersetup{pdftitle={\@title},pdfauthor={\@author}}}
\renewcommand{\maketitle}{%
  \thispagestyle{plain}%
  \begin{center}
    {\fontsize{17}{19.5}\selectfont\bfseries\@title\par}
    \vspace{4pt}
    {\small\@author\par}
    \ifx\@date\@empty\else
      \vspace{1pt}{\footnotesize\color{orgray}\@date\par}
    \fi
  \end{center}
  \vspace{4pt}\hrule height 0.55pt\vspace{7pt}
}
\makeatother

\pagestyle{fancy}
\fancyhf{}
\fancyhead[L]{\scriptsize\color{orgray}\obshorttitle}
\fancyhead[R]{\scriptsize\color{orgray}OR-Bench candidate problem}
\fancyfoot[C]{\scriptsize\thepage}
\renewcommand{\headrulewidth}{0.35pt}
\renewcommand{\footrulewidth}{0pt}

\fancypagestyle{plain}{
  \fancyhf{}
  \fancyfoot[C]{\scriptsize\thepage}
  \renewcommand{\headrulewidth}{0pt}
}

\renewenvironment{abstract}{%
  \noindent\begin{minipage}{\linewidth}\small
  \textbf{Abstract.}\hspace{0.35em}%
}{%
  \end{minipage}\par\vspace{5pt}
}

\newcommand{\keywords}[1]{%
  \noindent{\small\textbf{Keywords.}\hspace{0.35em}#1}\par\vspace{5pt}
}

\newenvironment{businessbrief}{%
  \par\vspace{4pt}\noindent
  \begin{minipage}{\linewidth}
  \hrule height 0.4pt\vspace{4pt}
  \small\textbf{Business-level description.}\hspace{0.35em}\itshape
}{%
  \par\vspace{4pt}\hrule height 0.4pt
  \end{minipage}\par\vspace{4pt}
}

\newcommand{\evidencebox}[1]{%
  \par\vspace{3pt}\noindent
  \colorbox{orlight}{%
    \parbox{\dimexpr\linewidth-2\fboxsep\relax}{%
      \small\textbf{Illustrative instance.}\hspace{0.35em}#1%
    }%
  }\par\vspace{4pt}
}

\newcolumntype{Y}{>{\raggedright\arraybackslash}X}
\newcolumntype{R}{>{\raggedleft\arraybackslash}X}

\AtBeginDocument{%
  \setlength{\abovedisplayskip}{6pt plus 2pt minus 2pt}
  \setlength{\belowdisplayskip}{6pt plus 2pt minus 2pt}
  \setlength{\abovedisplayshortskip}{3pt plus 1pt}
  \setlength{\belowdisplayshortskip}{4pt plus 1pt}
  \setlength{\jot}{2pt}
}

\title{Pair-State Consistency in Paired Combinatorial Logit Assortment Planning}
\author{Zhengzhong Ricky You}
\date{OR-Bench candidate problem, August 2026}
\shorttitle{Pair-State Consistency in Assortment Planning}

\begin{document}
\maketitle

\begin{abstract}
Assortment planning under the paired combinatorial logit (PCL) model couples a display decision with demand substitution across product pairs. When one member of a pair is offered and the other is not, the pair remains active and contributes the attraction weight of the offered product. We examine two incomplete formulations that replace the full pair-state system with either one global singleton term per offered product or terms for jointly offered pairs only. The PCL choice system of \citet{koppelman2000} and the assortment formulation of \citet{zhang2020} provide the reference model. On a six-product instance with a display limit of four, the unique optimal assortment is $\{3,5,6\}$ with expected margin 0.323647. Both incomplete formulations select $\{2,3,5,6\}$. Under the reference model, that assortment has expected margin 0.285975, which is 11.64\% below optimum. Omitting pair states changes both the choice probabilities and the resulting assortment decision.
\end{abstract}

\keywords{assortment optimization; paired combinatorial logit; discrete choice; revenue management; mixed-integer linear programming}

\section{Problem Motivation}

\begin{businessbrief}
A retailer must choose up to four products for a limited display. The demand science team has estimated an attraction weight and a normalized contribution margin for each product. It has also calibrated a paired combinatorial logit choice model with a dissimilarity parameter for every unordered product pair and an outside-option weight for customers who leave without purchasing. The retailer selects the assortment that maximizes expected contribution margin while accounting for substitution among the displayed products.
\end{businessbrief}

A display limit requires the retailer to compare discrete assortments in which the offered products compete for demand. Ranking these assortments therefore requires a choice model that maps each offered set to the corresponding purchase probabilities. The paired combinatorial logit (PCL) model of \citet{koppelman2000} provides such a representation by assigning a separate choice nest to each unordered pair of products. Using this pairwise representation, \citet{zhang2020} study assortment optimization under the PCL model. We adopt the standard PCL choice system and an exact linear-fractional reformulation as the reference model.

The pairwise nests introduce a state-consistency requirement. For each product pair, the contribution to demand depends on whether only the first product, only the second product, or both products are offered. The selected assortment must therefore determine these three states consistently. A formulation that omits a state or merges distinct states changes the underlying choice probabilities and may consequently select an assortment that is not optimal under the intended PCL model.

\section{Problem Definition}

Let $n\ge2$ be the number of products, let $N=\{1,\ldots,n\}$ be the product set, and let $\mathcal P=\{(i,j)\in N\times N\mid i<j\}$ be the set of unordered product pairs. The retailer may offer at most $K$ products, where $1\le K\le n$. Let $x=(x_i)_{i\in N}$ be the offer vector. Binary variable $x_i$ equals one when product $i\in N$ is offered and zero otherwise, so a feasible assortment satisfies $\sum_{i\in N}x_i\le K$.

Product $i\in N$ has attraction weight $v_i>0$ and normalized contribution margin $r_i\ge0$. The outside option has attraction weight $v_0>0$. Pair $(i,j)\in\mathcal P$ has dissimilarity parameter $\gamma_{ij}\in(0,1]$. We use one choice nest for each pair in $\mathcal P$, which fixes the scale of $v_0$.

For any feasible assortment $x$, the PCL model determines the purchase probabilities of the offered products and the outside option. The retailer chooses $x$ to maximize the resulting expected contribution margin.

\section{Reference Formulation}

Following \citet{koppelman2000} and \citet{zhang2020}, define the inclusive value of pair $(i,j)\in\mathcal P$ as
\begin{equation}
I_{ij}
=
\left(
v_i^{1/\gamma_{ij}}+v_j^{1/\gamma_{ij}}
\right)^{\gamma_{ij}}.
\label{eq:nest}
\end{equation}
When both products are offered, their contributions to the pair nest are
\begin{align}
\mu_{ij}^i
&=
I_{ij}
\frac{v_i^{1/\gamma_{ij}}}
{v_i^{1/\gamma_{ij}}+v_j^{1/\gamma_{ij}}},
\label{eq:massI}\\
\mu_{ij}^j
&=
I_{ij}
\frac{v_j^{1/\gamma_{ij}}}
{v_i^{1/\gamma_{ij}}+v_j^{1/\gamma_{ij}}}.
\label{eq:massJ}
\end{align}

Let $S=\{10,01,11\}$ be the set of active pair-state labels, where each label records the offer status $(x_i,x_j)$. For each $(i,j)\in\mathcal P$, binary variables $y_{ij}^{10}$, $y_{ij}^{01}$, and $y_{ij}^{11}$ indicate that only product $i$, only product $j$, or both products are offered. When neither product is offered, all three state variables equal zero.

For state $s\in S$, let $b_{ij}^s$ and $a_{ij}^s$ denote the denominator and expected-margin numerator coefficients. The denominator coefficients are $b_{ij}^{10}=v_i$, $b_{ij}^{01}=v_j$, and $b_{ij}^{11}=I_{ij}$. The numerator coefficients are $a_{ij}^{10}=r_iv_i$, $a_{ij}^{01}=r_jv_j$, and $a_{ij}^{11}=r_i\mu_{ij}^i+r_j\mu_{ij}^j$.
Because the state indicators are uniquely determined by $x$, the PCL denominator and expected-margin numerator are
\begin{align}
D(x)
&=
v_0+
\sum_{(i,j)\in\mathcal P}
\sum_{s\in S}
b_{ij}^s y_{ij}^s,
\label{eq:denom}\\
R(x)
&=
\sum_{(i,j)\in\mathcal P}
\sum_{s\in S}
a_{ij}^s y_{ij}^s.
\label{eq:numer}
\end{align}
The expected contribution margin is $R(x)/D(x)$. An exactly-one-offered pair contributes $v_i$ or $v_j$ to Eq.~\eqref{eq:denom} and the corresponding margin-weighted term to Eq.~\eqref{eq:numer}. A pair vanishes only when neither product is offered.

We apply the Charnes--Cooper transformation of \citet{charnes1962}. Let $t=1/D(x)$ and $w_{ij}^s=t y_{ij}^s$ for every $(i,j)\in\mathcal P$ and $s\in S$. Define
\begin{equation}
\overline D
=
v_0+
\sum_{(i,j)\in\mathcal P}
\max_{s\in S} b_{ij}^s,
\qquad
\underline t=\frac{1}{\overline D},
\qquad
\overline t=\frac{1}{v_0}.
\label{eq:tbounds}
\end{equation}
The inequalities $v_0\le D(x)\le\overline D$ imply $0<\underline t\le t\le\overline t$. The resulting mixed-integer linear program is
\begin{align}
\max\quad
&
\sum_{(i,j)\in\mathcal P}
\sum_{s\in S}
a_{ij}^s w_{ij}^s,
\label{eq:milpobj}\\
\text{subject to}\quad
&
v_0t+
\sum_{(i,j)\in\mathcal P}
\sum_{s\in S}
b_{ij}^s w_{ij}^s
=
1,
\label{eq:normalize}\\
&
\sum_{i\in N}x_i\le K,
\label{eq:cardinality}\\
&
y_{ij}^{11}\le x_i,\qquad
y_{ij}^{11}\le x_j,\qquad
y_{ij}^{11}\ge x_i+x_j-1
&&
\forall (i,j)\in\mathcal P,
\label{eq:and}\\
&
y_{ij}^{10}=x_i-y_{ij}^{11},\qquad
y_{ij}^{01}=x_j-y_{ij}^{11}
&&
\forall (i,j)\in\mathcal P,
\label{eq:states}\\
&
\underline t y_{ij}^s
\le
w_{ij}^s
\le
\overline t y_{ij}^s
&&
\forall (i,j)\in\mathcal P,\ s\in S,
\label{eq:scaleone}\\
&
t-\overline t(1-y_{ij}^s)
\le
w_{ij}^s
\le
t-\underline t(1-y_{ij}^s)
&&
\forall (i,j)\in\mathcal P,\ s\in S,
\label{eq:scaletwo}\\
&
x_i\in\{0,1\}
&&
\forall i\in N,
\label{eq:domainX}\\
&
y_{ij}^s\in\{0,1\},\qquad
w_{ij}^s\ge0
&&
\forall (i,j)\in\mathcal P,\ s\in S,
\label{eq:domainY}\\
&
\underline t\le t\le\overline t.
\label{eq:domainT}
\end{align}

Constraints~\eqref{eq:and} and \eqref{eq:states} select the unique state of every active pair. Constraints~\eqref{eq:scaleone} and \eqref{eq:scaletwo} impose $w_{ij}^s=t y_{ij}^s$. The normalization in Eq.~\eqref{eq:normalize} converts the fractional objective into the linear objective in Eq.~\eqref{eq:milpobj}. The resulting variables describe one assortment and its associated PCL choice probabilities.

\section{Failure of Incomplete Pair States}

Two incomplete formulations replace the three-state system with restricted pair-state representations. Using superscript $\mathrm{SP}$ for singleton plus pair, the denominator $D^{\mathrm{SP}}(x)$ and expected-margin numerator $R^{\mathrm{SP}}(x)$ of the first formulation are
\begin{align}
D^{\mathrm{SP}}(x)
&=
v_0+
\sum_{i\in N}v_i x_i+
\sum_{(i,j)\in\mathcal P}I_{ij}x_ix_j,
\label{eq:spDenom}\\
R^{\mathrm{SP}}(x)
&=
\sum_{i\in N}r_iv_i x_i+
\sum_{(i,j)\in\mathcal P}a_{ij}^{11}x_ix_j.
\label{eq:spNumer}
\end{align}
The global singleton term counts each offered product once, while the reference model contributes one exactly-one-offered term for every absent counterpart.

Using superscript $\mathrm{JP}$ for jointly offered pairs only, the denominator $D^{\mathrm{JP}}(x)$ and expected-margin numerator $R^{\mathrm{JP}}(x)$ of the second formulation are
\begin{align}
D^{\mathrm{JP}}(x)
&=
v_0+
\sum_{(i,j)\in\mathcal P}I_{ij}x_ix_j,
\label{eq:jpDenom}\\
R^{\mathrm{JP}}(x)
&=
\sum_{(i,j)\in\mathcal P}a_{ij}^{11}x_ix_j.
\label{eq:jpNumer}
\end{align}
The first and second incomplete formulations maximize $R^{\mathrm{SP}}(x)/D^{\mathrm{SP}}(x)$ and $R^{\mathrm{JP}}(x)/D^{\mathrm{JP}}(x)$, respectively, over the same feasible assortment set. Every pair with exactly one offered member disappears from the second choice system. Both formulations therefore change the numerator and denominator used to rank assortments.

The six-product instance is derived from the computational artifact associated with exact PCL assortment optimization under general linear constraints \citep{feng2026,artifact}. Product indices are remapped to $N=\{1,\ldots,6\}$, the display limit is $K=4$, and $v_0=3.845384$. The outside-option probability equals 0.25 when all six products are offered. All computations use the full-precision instance values. Reported objectives and absolute regret are rounded to six decimal places, and relative regret is rounded to two decimal places. Table~\ref{tab:decisions} compares the exact reference solution with the decisions of the two incomplete formulations.

\begin{center}
\captionof{table}{Assortments selected under the reference and incomplete formulations.}
\label{tab:decisions}
\begin{tabularx}{0.94\linewidth}{lY r}
\toprule
Formulation & Selected products & Objective under the formulation \\
\midrule
Reference PCL formulation & $\{3,5,6\}$ & 0.323647 \\
Singleton-plus-pair formulation & $\{2,3,5,6\}$ & 0.245658 \\
Jointly-offered-pairs-only formulation & $\{2,3,5,6\}$ & 0.188962 \\
\bottomrule
\end{tabularx}
\end{center}

Each incomplete formulation is solved to optimality under its own objective. The two altered choice systems select the same four-product assortment.

\evidencebox{The exact mixed-integer linear program and complete enumeration of all 57 feasible assortments identify the unique optimum $\{3,5,6\}$ with expected margin 0.323647. Both incomplete formulations select $\{2,3,5,6\}$. Under the reference PCL model, that assortment has expected margin 0.285975, giving absolute regret 0.037672 and relative regret 11.64\%.}

The display-capacity constraint holds for all three decisions. The discrepancy follows from the choice probabilities assigned to exactly-one-offered pairs. Preserving the complete pair-state system is therefore necessary for the optimization model to evaluate the calibrated PCL demand model.

Supporting materials for the case study, including the derivative data, source document, reference solver, and verification instructions, are available in the \href{https://github.com/Zhengzhong-You/OR-Bench/tree/main/problems/02_pcl_assortment}{corresponding folder of the OR-Bench GitHub repository}.

\begin{thebibliography}{9}\small

\bibitem[Charnes and Cooper(1962)]{charnes1962}
Charnes A, Cooper WW (1962)
Programming with linear fractional functionals.
\emph{Naval Research Logistics Quarterly} 9(3--4), 181--186.

\bibitem[Feng et al.(2026a)]{feng2026}
Feng J, Che A, Chen ZL (2026a)
Assortment optimization with general linear constraints under the paired combinatorial logit choice model.
\emph{INFORMS Journal on Computing}, Articles in Advance.

\bibitem[Feng et al.(2026b)]{artifact}
Feng J, Che A, Chen ZL (2026b)
Code and data repository for assortment optimization with general linear constraints under the paired combinatorial logit choice model.
\emph{INFORMS Journal on Computing Software and Data Repository}.

\bibitem[Koppelman and Wen(2000)]{koppelman2000}
Koppelman FS, Wen CH (2000)
The paired combinatorial logit model: Properties, estimation and application.
\emph{Transportation Research Part B: Methodological} 34(2), 75--89.

\bibitem[Zhang et al.(2020)]{zhang2020}
Zhang H, Rusmevichientong P, Topaloglu H (2020)
Assortment optimization under the paired combinatorial logit model.
\emph{Operations Research} 68(3), 741--761.

\end{thebibliography}

\end{document}
